\section{Axial variables and the comparison family}\label{a:section}

We first encode the fields and the components removed by polarization in
three scalar potentials.  The metric and second fundamental form in this
section are the original data until a subscript~$0$ is introduced.
Throughout, $\eta$ is the prescribed generator of period $2\pi$.

\subsection{The orbit space and its smooth angular coordinate}

\begin{lemma}[Orbit topology]\label{a:topology}
The quotient of $\Omega$ by the circle action is a half annulus: its
boundary consists of the compact image of $S$ and two noncompact axis
faces, each joining an endpoint of that image to infinity.  There are no
exceptional orbits.  The regular part of $\Omega$ is a trivial principal
circle bundle over the interior of this quotient.  It admits a global
angular coordinate that agrees with the prescribed azimuth on a smaller
asymptotic end and with smooth polar gauges near the axes.
\end{lemma}

\begin{proof}
Truncate the given end at a sufficiently large rotation sphere.  Both
boundary spheres carry ordinary effective rotations.  Here is a direct
verification of the assertion for an invariant sphere.  Its Euler
characteristic forces fixed points.  At a fixed point the tangent
representation is a planar rotation with integer weight.  That weight
has absolute value one: otherwise a nonidentity group element would
fix the point and have identity differential both on the sphere and on
its invariant inward normal, hence would be an isometry with identity
first jet and therefore the identity on the connected ambient
manifold.  At a nonfixed point a finite stabilizer acts trivially on the
orbit tangent.  It acts trivially also on the transverse tangent line,
by orientation, and on the inward normal.  The same first-jet argument
excludes such a stabilizer.  The slice theorem now makes the quotient
of the sphere an interval, with an ordinary rotation disk at either
endpoint; the free circle bundle over its interior is a product.
Consequently the action extends over an ordinary rotation ball.

Cap the two boundary spheres equivariantly, using invariant collars.
The resulting closed three-manifold is simply connected.  Its orbit
space is a compact orientable surface with nonempty boundary, the
boundary arising from fixed axes.  The only other possible singular
orbits are isolated exceptional fibers, represented by interior cone
points of finite cyclic order.  The orbifold fundamental group of this
surface is a quotient of the fundamental group of the three-manifold.
Indeed, away from axes and exceptional fibers, loops lift to paths;
an endpoint displacement in a fiber can be closed within that fiber
without changing the projected loop.  Filling an axis tube introduces
no nontrivial projected fiber relation.  Filling an exceptional solid
torus imposes the relation that the corresponding meridian raised to
the isotropy order is trivial, as seen in its cyclic disk cover.
Van Kampen therefore gives the asserted surjection onto the orbifold
group.

An orientable surface with nonempty boundary and cyclic cone points
has orbifold group equal to the free product of its ordinary free
surface group and the cyclic groups at those points.  Triviality of
this group implies that the surface is a disk and that there are no
cone points.  Removing the two cap quotients gives the half-annulus
description.  Its regular interior is contractible, so its principal
circle bundle is trivial.

For the last smoothness assertion, use equivariant normal disks near
each axis and a signed transverse polar radius there.  Changes of
angular reference are smooth functions on the quotient, even in that
signed radius.  Their local phase lifts can be patched by a partition
of unity; the only obstruction is the integral circle cocycle, which
vanishes on this orbit rectangle.  The same construction relative to
a far-end strip preserves the given azimuth there.  Thus the product
gauge has the stated compatibility with the smooth axes.  No
isothermal, Weyl, or other metric coordinates have been assumed.
\end{proof}

Off the axis write, with horizontal lifts understood,
\begin{equation}\label{a:metric-splitting}
 g=h+X^2(\dd\phi+\mathcal A)^2,
 \qquad X=|\eta|_g,\qquad e=X^{-1}\eta.
\end{equation}
Orient the meridional metric $h$ by
$\dd V_g=\dd A_h\wedge X(\dd\phi+\mathcal A)$.
Thus $*_h$ sends the first covector of a positive orthonormal
meridional coframe to the second.  A subscript $H$ denotes a
horizontal component, and horizontal vector fields will be identified
with their $h$-dual covectors when applying $*_h$.

\subsection{The conserved potentials and their axis behavior}

Perform the constant duality rotation in Theorem~\ref{m:main} and
write $E,B$ for the resulting fields in this section.  Their charges
are $Q,0$.  The rotation preserves $|E|^2+|B|^2$ and $E\times B$,
so it preserves the electrovacuum constraints.  It gives
$B=O_3(r^{-3})$ on the end.

\begin{lemma}[Magnetic primitive and total angular momentum]
\label{a:angular-momentum}
There is a smooth invariant one-form $\alpha$ with
\[
 \dd\alpha=*_g B^\flat,\qquad \alpha=O_2(r^{-2}).
\]
The flux
\begin{equation}\label{a:total-flux}
 \frac1{8\pi}\int_\Sigma(K-\tau g)(\nu,\eta)\,\dd A_g
 +\frac1{4\pi}\int_\Sigma\alpha(\eta)g(E,\nu)\,\dd A_g
\end{equation}
is independent of the invariant enclosing surface $\Sigma$ and of
the smooth invariant gauge for $\alpha$.  It equals $J_{\rm ADM}$
at infinity, and hence equals the prescribed $J$.
\end{lemma}

\begin{proof}
The magnetic flux form is closed by $\divg_gB=0$.  Since
$\Omega$ retracts onto a sphere, its only two-dimensional de Rham
period is the magnetic flux, which is zero after rotation.  The form
is therefore exact.  The decay can be imposed without sacrificing
smoothness.  On a rescaled annulus, radial homotopy reduces the
primitive problem to a two-form on a fixed sphere; its zero integral
is exactly the condition for a spherical primitive.  These homotopy
operators, followed by a fixed spherical elliptic inverse, have
uniform bounds through the required derivatives.  Scaling back a
two-form of Cartesian order $r^{-3}$ gives a primitive of order
$r^{-2}$, with its first two derivatives of the corresponding orders.
On consecutive overlapping annuli, the difference of the primitives
is a closed one-form and hence is the differential of a function.
After subtracting a constant, its rescaled derivatives are bounded.
Use cutoff multiples of these functions on the overlaps to glue the
annular primitives.  The annuli can be chosen with bounded overlap,
so these corrections retain the bounds.  The same argument glues the
end primitive to any interior primitive.  All ingredients are smooth
on each finite annulus; this avoids any demand for uniform bounds on
unspecified higher derivatives at infinity.  Averaging over the
circle makes $\alpha$ invariant and preserves its equation and decay.

Set $\chi=\alpha(\eta)$.  Cartan's formula gives
\[
 \dd\chi=-\iota_\eta(*_gB^\flat),
 \qquad (E\times B)\cdot\eta=-E\cdot\nabla\chi.
\]
The contraction of the momentum constraint with the Killing field,
together with $\divg E=0$, now gives
\[
 \divg_g\big((K-\tau g)(\eta,\cdot)^\sharp+2\chi E\big)=0.
\]
The divergence theorem proves conservation of
Equation~\eqref{a:total-flux}.  If $\alpha$ is changed by a closed
invariant one-form $\beta$, then
$\dd(\beta(\eta))=-\iota_\eta\dd\beta=0$.
The constant $\beta(\eta)$ is zero on the axis, and thus everywhere.
Finally $\alpha(\eta)=O(r^{-1})$, while $E\cdot\nu=O(r^{-2})$;
the electromagnetic surface term is $O(r^{-1})$ on the end spheres.
This proves its disappearance in the limiting flux.  The proof does
not identify the bare gravitational flux on $S$ with the flux at
infinity.
\end{proof}

Let $l=K(e,\cdot)|_H$.  The global potentials are defined by
\begin{equation}\label{a:potentials}
 \dd\psi=X*_hE_H,\qquad
 \dd\chi=X*_hB_H,\qquad
 \omega:=\dd z+\psi\,\dd\chi-\chi\,\dd\psi=X^2*_h l.
\end{equation}
For any positive axial radius $X$ define
\begin{equation}\label{a:target}
 \begin{gathered}
 \mathcal D_X=(X^{-1}\dd\psi,X^{-1}\dd\chi,X^{-2}\omega),
 \qquad u=\log X,\\
 G=\dd u^2+e^{-2u}(\dd\psi^2+\dd\chi^2)
       +e^{-4u}(\dd z+\psi\dd\chi-\chi\dd\psi)^2.
 \end{gathered}
\end{equation}
The last expression is a Riemannian metric on $\R^4$, with
coordinates $(u,\psi,\chi,z)$.

We verify global integrability and the normalizations in
Equation~\eqref{a:potentials}.  For an invariant field the divergence
equation is equivalent to the closedness of $X*_hE_H$, and similarly
for $B$.  The second of these forms is $\dd(\alpha(\eta))$, by
Cartan's identity with the stated orientation.  The axial momentum
equation gives
\[
 \dd(X^2*_h l)=2\,\dd\psi\wedge\dd\chi,
 \qquad
 \dd\big(X^2*_h l+2\chi\dd\psi\big)=0.
\]
Consequently
$X^2*_h l+\chi\dd\psi-\psi\dd\chi$ is closed.  The
meridional rectangle is simply connected, so all three defining
equations integrate there.

Orient a crossing arc so that its tangent is the metric dual of
$*_h\nu^\flat$, where $\nu$ is the horizontal normal pointing
toward infinity.  Rotating the arc gives area element
$2\pi X\,\dd s_h$.  The electric flux divided by $4\pi$ is
therefore $\frac12\int\dd\psi$, and the total angular flux is
\[
 \frac14\int(\omega+2\chi\dd\psi)
       =\frac14\int\dd(z+\chi\psi).
\]
All potentials are constant on each axis face.  Order those faces
by this crossing orientation.  Since $\chi=\alpha(\eta)=0$ on
both axes, additive constants can and will be chosen so that their
values are
\begin{equation}\label{a:axis-values}
 (\psi,\chi,z)=(-Q,0,-2J)\quad\hbox{and}\quad(Q,0,2J).
\end{equation}
The same computation on end-coordinate spheres proves that these
are the original flux normalizations.

\begin{lemma}[Axis and end regularity]\label{a:axis-regularity}
In a compact ordinary axis tube, let $y$ be signed transverse polar
radius and let $x$ be a coordinate along the axis.  The functions
\[
 \psi-\psi_{\rm ax},\qquad\chi-\chi_{\rm ax}
\]
are smooth even functions of $y$ of order $y^2$, while
\[
 z+\psi_{\rm ax}\chi-\chi_{\rm ax}\psi-z_{\rm ax}
\]
is smooth even of order $y^4$.  These statements hold with smooth
longitudinal derivatives and at the boundary poles in one-sided
charts.  On the asymptotic end, the potentials are bounded, and near
either pole their differences from the adjacent axis values obey
\begin{equation}\label{a:end-regularity}
 \psi-\psi_{\rm ax}=O(\sin^2\vartheta),\qquad
 \chi=O(r^{-1}\sin^2\vartheta),\qquad
 z-z_{\rm ax}=O(\sin^2\vartheta).
\end{equation}
The corrected central coordinate satisfies the stronger uniform estimate
\begin{equation}\label{a:end-central}
 \widetilde z:=z+\psi_{\rm ax}\chi-\chi_{\rm ax}\psi-z_{\rm ax}
       =O(\sin^4\vartheta).
\end{equation}
The estimates have the differentiated versions obtained from
Equation~\eqref{a:potentials}; in particular their rescalings to end
annuli have uniform bounds through the orders used below.
\end{lemma}

\begin{proof}
An invariant smooth scalar is even in $y$.  A smooth invariant
horizontal vector field has an odd radial component and an even
longitudinal component.  Since $X/y$ is smooth positive even,
$X*_hE_H$ and $X*_hB_H$ have radial component $O(y)$ and
longitudinal component $O(y^2)$ with the corresponding parities.
Integration from the axis gives the first two assertions.

Smoothness of an invariant symmetric tensor gives an angular--axis
component $l_x=O(y)$ and an angular--radial component $l_y=O(y^2)$.
After multiplication by $X^2$ and rotation by $*_h$, $\omega$ has
radial component $O(y^3)$ and longitudinal component $O(y^4)$.
Put $\delta\psi=\psi-\psi_{\rm ax}$ and
$\delta\chi=\chi-\chi_{\rm ax}$.  Then
\[
 \dd(z+\psi_{\rm ax}\chi-\chi_{\rm ax}\psi)
       =\omega-\delta\psi\,\dd\chi
                    +\delta\chi\,\dd\psi.
\]
The last two terms have the same radial and longitudinal orders as
$\omega$.  Integration proves the order-four assertion.

For the end estimates, perform the same argument in rescaled
annuli.  The leading electric field is the radial Coulomb field, so
the leading angular derivative of $\psi$ is
$Q\sin\vartheta\,\dd\vartheta$.  The remainder contributes a
uniformly bounded function of order $\sin^2\vartheta$ at a pole.
The decaying invariant magnetic primitive gives the sharper estimate
for $\chi$: its normalized angular component vanishes linearly in
transverse radius, and $\chi=\alpha(\eta)$ vanishes quadratically.
Finally $X^2l=O(r^{-1}\sin^2\vartheta)$ as a horizontal
covector, with the stronger radial and longitudinal polar orders
just proved.  Integration along a meridian of a coordinate sphere
therefore bounds $z$ and gives its order-two polar difference.
For the stronger corrected estimate, fix a dyadic annulus of radius
$R$ and use coordinates scaled by $R$, with signed transverse coordinate
$y$.  The rescaled Cartesian tensor $R^3K(R\,\cdot)$ has uniform $C^3$
bounds and the same invariant tensor parity.  The pullback of
$\omega=X^2*_hl$ therefore has transverse component $O(y^3)$ and
longitudinal component $O(y^4)$, uniformly in $R$: the factors from
$X^2$ and a covector pullback cancel the overall $R^{-3}$ decay of $K$.
The first two potential estimates and their differentiated versions give
these same orders for
$-(\psi-\psi_{\rm ax})\dd\chi
 +(\chi-\chi_{\rm ax})\dd\psi$.
Integrating their sum with $\omega$ from $y=0$ gives
$\widetilde z=O(y^4)$ uniformly.  On the unit-scale annulus
$y$ is comparable to $\sin\vartheta$, proving
Equation~\eqref{a:end-central}.  Smooth polar coordinate changes with
uniformly bounded transverse ratios preserve these estimates.
Differentiating the defining equations supplies the radial decay
and angular derivative bounds.  The assumptions $g-\delta=O_4(r^{-1})$,
$K=O_3(r^{-3})$ and $E,B$ with three controlled derivatives give
all the asserted rescaled bounds.
\end{proof}

\subsection{Polarization and its exact source terms}

Define
\begin{equation}\label{a:polarized-data}
 g_0=h+X^2\dd\phi^2,
 \qquad
 k_0=K|_{H\times H}+K(e,e)X^2\dd\phi^2.
\end{equation}
Thus $\mathcal A$ and $l$ are removed from the metric and tensor,
but retained in the map and in the velocity variables below.

\begin{proposition}[Polarization]\label{a:polarization}
The data $(g_0,k_0)$ extend smoothly through the axis, have the
stated asymptotic decay, and have the same ADM energy as $(g,K)$.
For every invariant surface, its area and future expansion are
unchanged.  Let $\mu_0,J_0$ be the constraint densities of
$(g_0,k_0)$, with $J_0$ a covector.  Write
$\dd\mathcal A=f_a\dd A_h$ and set
$s=(-B_e,E_e,-Xf_a/2)$.  Then
\begin{equation}\label{a:constraints}
 16\pi\mu_0=2\big(|\mathcal D_X|_{g_0}^2+|s|^2\big),
 \qquad
 8\pi J_0=2s\cdot\mathcal D_X.
\end{equation}
The source combinations on the right are smooth across the axis.
On a signed axis tube $s_1,s_2$ are odd of order $y$, and $s_3$
is even of order $y^2$.  In particular the polarized data satisfy
the dominant energy condition.
\end{proposition}

\begin{proof}
First consider smoothness.  In the normalized polar coframe
$(\dd y,\dd x,y\dd\phi)$, the angular--radial metric coefficient
is even of order $y^2$, and the angular--longitudinal coefficient
is odd of order $y$.  The ratio $X/y$ is smooth positive even.
Taking the horizontal Schur complement of the angular coefficient
therefore preserves all polar parities and the equality of the
radial and angular diagonal coefficients at the axis.  The same
calculation for the tensor gives the smooth lift of $k_0$.
These statements apply uniformly on rescaled end annuli.  Apparent
factors $y^{-1}$ in Cartesian differentiation are harmless: the
Taylor terms permitted by the polar parities are smooth transverse
Cartesian polynomials, and Taylor remainder estimates bound every
remaining divided derivative.  This also proves the required
differentiated end decay.  In particular
\[
 \mathcal A_y=O(yr^{-3}),\qquad
 \mathcal A_x=O(r^{-2}),\qquad Xf_a=O(r^{-2})
\]
on the end.  On compact axis tubes $\mathcal A_y$ is smooth odd
and $\mathcal A_x$ smooth even, so $f_a$ is odd of order $y$.
Smooth invariant vector fields have angular component odd of
order $y$, proving the stated parities of $s$.

An invariant surface is generated by a meridional curve, and its
area is $2\pi$ times that curve's length in $X^2h$.  This is the
same for $g$ and $g_0$.  First variation of this expression for
arbitrary invariant normal speeds shows that its pointwise mean
curvature is also unchanged.  Its tangent space is generated by
one horizontal unit vector and $e$, so its tangential trace uses
only $K|_{H\times H}$ and $K(e,e)$.  This proves preservation of
future expansion.

For the energy, use the end azimuth fixed in
Lemma~\ref{a:topology}.  Let $\mathcal R$ be azimuthal reflection.
The difference of $g$ from its reflection average is precisely
the angular cross tensor, which is odd under $\mathcal R$.
The linear ADM flux on a coordinate sphere is invariant under
this Euclidean reflection, and hence is zero on that odd tensor.
The difference of the reflection average from $g_0$ is
$X^2\mathcal A\otimes\mathcal A$ in horizontal coordinate
components.  It has Cartesian order $O_2(r^{-2})$, since it is
quadratic in the normalized cross coefficients $O(r^{-1})$.
Its ADM flux is $O(r^{-1})$.  Thus the energies agree.  The polar
estimates above make the calculation uniform at both end poles.

It remains to prove the source identities, including their signs.
The horizontal Koszul formulas, for invariant horizontal fields
$a,b$, are
\begin{align*}
 (\nabla_a b)^e&=-\tfrac12X\dd\mathcal A(a,b),&
 (\nabla_a e)^H=(\nabla_e a)^H
       &=\tfrac12X\dd\mathcal A(a,\cdot)^\sharp,&
 \nabla_e e&=-\nabla_h\log X.
\end{align*}
Together with the base connection these give
\begin{equation}\label{a:deletion-identities}
 \begin{gathered}
 R_{g_0}-R_g=\tfrac12X^2f_a^2,\qquad
 |K|_g^2-|k_0|_{g_0}^2=2|l|_h^2,\qquad
 \tr_{g_0}k_0=\tr_gK,\\
 8\pi(J_0)_i=8\pi(J_{\rm con})_i
                       +X(\dd\mathcal A)_{ij}l^j
                       \quad(i\text{ horizontal}),\qquad
 J_0(e)=0.
 \end{gathered}
\end{equation}
For example the last vanishing also follows from the azimuthal
reflection symmetry of $(g_0,k_0)$.  The scalar constraint and
\begin{align*}
 |\mathcal D_X|^2&=|E_H|^2+|B_H|^2+|l|^2,\\
 (E\times B)_H&=E_e*_hB_H-B_e*_hE_H,\qquad
 (\dd\mathcal A)_{ij}l^j=-f_a(*_hl)_i
\end{align*}
now give Equation~\eqref{a:constraints} exactly.

For clarity, the normalized first two rows of $\mathcal D_X$
need not individually be smooth three-dimensional covectors at
the axis.  Their radial coefficients are even and their
longitudinal coefficients odd.  Their squared norms, and their
products with the corresponding odd $s_i$, are smooth invariant
scalars and covectors.  The third row is a smooth horizontal
covector: its radial coefficient is odd of order $y$, and its
longitudinal coefficient is even of order $y^2$.  Thus all source
combinations in Equation~\eqref{a:constraints} are smooth.
Finally $2|s\cdot\mathcal D_X|\le |s|^2+|\mathcal D_X|^2$
implies $\mu_0\ge|J_0|_{g_0}$.
\end{proof}

\subsection{The spacetime quotient and local reconstruction}

We record the corresponding spacetime equations with all signs
fixed.  This is a local assertion on regular orbits and does not
assume a stationary development of the initial data.  Write
\[
 \overline g=H+X^2(\dd\phi+\mathbf A)^2,
 \qquad \gamma=X^2H,
\]
where $H$ is a Lorentz metric of dimension three and signature
$(-++)$.  Its volume form is ordered by a future time covector,
then the oriented horizontal spatial coframe.  The Maxwell field
in this subsection has already undergone the constant duality
rotation.

\begin{proposition}[Einstein--wave-map quotient]
\label{a:quotient-proposition}
For an invariant electrovacuum spacetime the axial potentials
satisfy
\begin{equation}\label{a:quotient}
 \dd\chi=-\iota_\eta F,\qquad
 \dd\psi=\iota_\eta *_4F,\qquad
 \omega=-\tfrac12X^3*_H\dd\mathbf A,
 \qquad \mathcal D_X(n)=s,
\end{equation}
on every invariant spacelike slice, where $n$ is its horizontal
future unit normal.  With $\Phi=(u,\psi,\chi,z)$, the equations
reduce to
\begin{equation}\label{a:wave}
 \Ric_\gamma=2\Phi^*G,
 \qquad \Box_\gamma\Phi=0.
\end{equation}
The second expression is the wave-map tension for the target
$G$.  Conversely, these reduced equations reconstruct locally
an invariant Einstein--Maxwell metric and field, up to angular
and electromagnetic gauge.  Their slice data recover the
original metric, tensor, and fields from the polarized data,
the potentials, and their normal velocities.
\end{proposition}

\begin{proof}
Let $(\theta^0,\theta^1,\theta^2)$ be an oriented Lorentz
orthonormal coframe with $\theta^0(n)=1$, and put
$\theta^3=X(\dd\phi+\mathbf A)$.  The relevant signs are
\[
 *_H\theta^0=-\theta^1\wedge\theta^2,\qquad
 *_H\theta^1=-\theta^0\wedge\theta^2,\qquad
 *_H\theta^2=\theta^0\wedge\theta^1.
\]
Set $p=X^{-1}\dd\chi$, $t=X^{-1}\dd\psi$, and
$v=X^{-2}\omega$.  Then
\begin{equation}\label{a:field-reconstruction}
 F=p\wedge\theta^3+*_Ht,\qquad
 *_4F=-t\wedge\theta^3+*_Hp,\qquad
 \dd\mathbf A=2X^{-1}*_Hv.
\end{equation}
Contraction with $n$ and restriction to its spatial slice give
exactly $E^\flat=\iota_nF$ and
$B^\flat=*_g(F|_{T\Omega})$.  In particular
$t(n)=-B_e$ and $p(n)=E_e$.
Writing
$\dd\mathbf A=\theta^0\wedge a+f_a\theta^1\wedge\theta^2$,
the mixed connection formula gives $l=Xa/2$ for the positive
future-normal convention.  Consequently
\[
 \omega|_{T\Omega}=X^2*_hl,
 \qquad X^{-2}\omega(n)=-Xf_a/2.
\]
This proves the restrictions in Equation~\eqref{a:quotient}.

For completeness, the integrability of the corrected twist in
spacetime follows from the mixed Einstein equation.  The mixed
Ricci component is
\[
 (\Ric_4)_{ie}
   =\frac1{2X^2}\nabla_H^j\big(X^3(\dd\mathbf A)_{ij}\big)
   =X^{-2}\nabla_H^j(*_H\omega)_{ij}.
\]
Einstein--Maxwell gives it as
$-2(*_Ht)_{ij}p^j$.  These identities are equivalent to
$\dd\omega=2\dd\psi\wedge\dd\chi$, and hence to the
closedness of
$\omega+\chi\dd\psi-\psi\dd\chi$.
Thus $z$ exists locally, with its constants fixed by the slice.

The horizontal parts of the Maxwell equations give
$\divg_Ht=2\langle p,v\rangle_H$ and
$\divg_Hp=-2\langle t,v\rangle_H$, while
$\dd^2\mathbf A=0$ gives $\divg_H(X^{-1}v)=0$.
For a covector $\beta$, the conformal change in dimension three
obeys $\divg_\gamma\beta=X^{-3}\divg_H(X\beta)$.
Therefore
\begin{align}\label{a:map-divergences}
 \divg_\gamma(X^{-4}\omega)&=0,\notag\\
 \divg_\gamma(X^{-2}\dd\psi)
       &=2X^{-4}\langle\dd\chi,\omega\rangle_\gamma,\\
 \divg_\gamma(X^{-2}\dd\chi)
       &=-2X^{-4}\langle\dd\psi,\omega\rangle_\gamma.\notag
\end{align}

The horizontal and fiber curvature contractions of the connection
formulas are
\begin{align}\label{a:ricci-blocks}
 (\Ric_4)_{ij}
   &=(\Ric_H)_{ij}-X^{-1}\nabla_i\nabla_jX
          -\tfrac12X^2(\dd\mathbf A)_{ik}(\dd\mathbf A)_j{}^k,
          \notag\\
 (\Ric_4)_{ee}
   &=-X^{-1}\Box_HX
          +\tfrac14X^2(\dd\mathbf A)_{ij}(\dd\mathbf A)^{ij}.
\end{align}
One can see the coefficients directly before contraction: the
horizontal sectional numerator has correction
$-3X^2\dd\mathbf A(a,b)^2/4$, and the horizontal--fiber
block is
$-X^{-1}\nabla_a\nabla_bX+
X^2(\dd\mathbf A)_{ak}(\dd\mathbf A)_b{}^k/4$.
The Lorentz Hodge identity
\[
 (*_Ht)_{ik}(*_Ht)_j{}^k=t_it_j-H_{ij}|t|_H^2
\]
shows that the two Einstein--Maxwell Ricci blocks in
Equation~\eqref{a:ricci-blocks} equal, respectively,
\[
 2(p_ip_j+t_it_j)-H_{ij}(|p|_H^2+|t|_H^2),
 \qquad |p|_H^2+|t|_H^2.
\]
Thus the fiber equation reads
$X^{-1}\Box_HX=-|p|_H^2-|t|_H^2-2|v|_H^2$.
Since $u=\log X$ and $\gamma=e^{2u}H$, this becomes
\begin{equation}\label{a:u-wave}
 \Box_\gamma u
  =-e^{-2u}(|\dd\psi|_\gamma^2+|\dd\chi|_\gamma^2)
       -2e^{-4u}|\omega|_\gamma^2.
\end{equation}
The conformal Ricci formula
\[
 \Ric_\gamma=\Ric_H-\Hess_Hu+\dd u\otimes\dd u
                -(\Box_Hu+|\dd u|_H^2)H
\]
then cancels the remaining scalar multiples of $H$ and gives
\[
 \Ric_\gamma=2(\dd u\otimes\dd u+p\otimes p+t\otimes t+v\otimes v)
             =2\Phi^*G.
\]
Equations~\eqref{a:map-divergences} and~\eqref{a:u-wave}
are exactly the four Euler--Lagrange equations of the target
metric in Equation~\eqref{a:target}.  This proves
Equation~\eqref{a:wave}.

Conversely, start with Equation~\eqref{a:wave}, put
$H=X^{-2}\gamma$, and define $p,t,v$ as above.  The first
equation in~\eqref{a:map-divergences} makes
$2X^{-1}*_Hv$ closed, so it has a local connection primitive
$\mathbf A$.  Define $F$ by
Equation~\eqref{a:field-reconstruction}.  The remaining divergence
equations and the identities
$\dd^2\psi=\dd^2\chi=0$ give both Maxwell equations.
The identity $\dd\omega=2\dd\psi\wedge\dd\chi$ supplies
the mixed Einstein equation.  Reversing the horizontal and fiber
calculations supplies the other Einstein equations.

The slice reconstruction is explicit.  Its horizontal tensor is
the second fundamental form of its slice in $H$; the other
components are
\begin{equation}\label{a:reconstruction}
 K(e,e)=n(\log X),\qquad
 X^2*_hK(e,\cdot)|_H=\omega|_{T\Omega},\qquad
 (X^{-1}\dd\psi(n),X^{-1}\dd\chi(n),X^{-2}\omega(n))=s.
\end{equation}
The spatial potential equations recover $E_H,B_H$, and the first
two velocities recover $B_e,E_e$.  The third recovers the spatial
connection curvature.  A chosen connection on a simply connected
meridian is then fixed up to angular gauge.  Finally the inverse
constant duality rotation recovers the original ordered pair of
electric and magnetic fields.
\end{proof}

\subsection{Target geometry}

\begin{lemma}\label{a:target-geometry}
The target $(\R^4,G)$ is complete, simply connected, and has
nonpositive curvature.  In particular any two of its points are
joined by a unique geodesic, and squared distance is jointly
convex along pairs of geodesics.
\end{lemma}

\begin{proof}
A path of bounded $G$-length has bounded $u$.  The same length
then controls the variations of $\psi$ and $\chi$.  Since these
coordinates are bounded, control of
$\dd z+\psi\dd\chi-\chi\dd\psi$ also controls the variation
of $z$.  A finite-length escaping path is therefore impossible,
which proves completeness.  The underlying manifold is $\R^4$.

Let $E_0,E_1,E_2,E_3$ be the frame dual to
$(\dd u,e^{-u}\dd\psi,e^{-u}\dd\chi,e^{-2u}\omega)$.
Its nonzero brackets are
\[
 [E_0,E_1]=E_1,\quad [E_0,E_2]=E_2,\quad
 [E_0,E_3]=2E_3,\quad [E_1,E_2]=-2E_3.
\]
In the bivector order $(01,02,12,03,13,23)$, Koszul's formula
gives the curvature operator
\[
 \begin{pmatrix}
 -1&0&0&0&0&-1\\
 0&-1&0&0&1&0\\
 0&0&-4&2&0&0\\
 0&0&2&-4&0&0\\
 0&1&0&0&-1&0\\
 -1&0&0&0&0&-1
 \end{pmatrix}.
\]
Its quadratic form on $(a,b,c,d,e,f)$ is
$-(a+f)^2-(b-e)^2-(c+d)^2-3(c-d)^2$, so it is
nonpositive.  The geodesic and convexity assertions follow from
the Cartan--Hadamard theorem and the second variation formula.
\end{proof}

\subsection{The Kerr--Newman comparison maps}

The models below are comparison objects.  Their stationarity is
not a hypothesis on the given data.  Set
\[
 C=Q^4+4J^2,\qquad
 F_*(R)=\frac12\sqrt{R^2+2Q^2+C/R^2}.
\]

\begin{lemma}[Model parameters]\label{a:model-parameters}
For every $R>0$ with $R^2>\sqrt C$, define
\[
 M=F_*(R),\qquad a=-J/M,\qquad
 b=\sqrt{M^2-a^2-Q^2}.
\]
Then
\begin{equation}\label{a:parameter-identities}
 b=\frac{R^4-C}{4MR^2}>0,\qquad
 (M+b)^2+a^2=R^2,\qquad F_*'(R)=\frac bR.
\end{equation}
The function $b(R)$ is strictly increasing and maps this branch
onto $(0,\infty)$.
\end{lemma}

\begin{proof}
Substitute $4M^2=R^2+2Q^2+C/R^2$ and $a^2=J^2/M^2$.
Clearing denominators gives the first two identities.  Ordinary
differentiation gives the third and also
\[
 b'(R)=
 \frac{R^8+4Q^2R^6+6CR^4+4Q^2CR^2+C^2}
 {2R^2(R^4+2Q^2R^2+C)^{3/2}}>0.
\]
At the lower endpoint $b\to0$ (also when $C=0$, where the
endpoint is $R=0$); as $R\to\infty$, $b\sim R/2$.
\end{proof}

Use the azimuth oriented by $\eta$ and the standard electrically
charged \KN\ solution with the parameter $a$ just fixed.  In
coordinates
\[
 r=M+b\cosh\sigma,\qquad \sigma>0,\qquad0<\vartheta<\pi,
\]
put
\begin{align*}
 \Delta&=r^2-2Mr+a^2+Q^2=b^2\sinh^2\sigma,\\
 P&=r^2+a^2\cos^2\vartheta,\\
 W&=(r^2+a^2)^2-a^2\Delta\sin^2\vartheta.
\end{align*}
Its metric and a Maxwell potential are
\begin{equation}\label{a:kn-metric}
 \begin{split}
 \overline g_*={}&-\frac{P\Delta}{W}\dd t^2
       +P(\dd\sigma^2+\dd\vartheta^2)
       +\frac WP\sin^2\vartheta
            \left(\dd\phi-\frac{a(2Mr-Q^2)}W\dd t\right)^2,\\
 \mathfrak a_*={}&\frac{Qr}{P}
                  (\dd t-a\sin^2\vartheta\,\dd\phi),
 \qquad \mathcal F_*=\dd\mathfrak a_*.
 \end{split}
\end{equation}
We use the classical explicit fact~\cite{NewmanEtAl1965} that
Equation~\eqref{a:kn-metric} is an Einstein--Maxwell solution.
For the field-equation verification and the Boyer--Lindquist form,
see also~\cite[Section~2.2 and Equation~(3.2)]{AdamoNewman2014},
with the opposite metric signature.
The reduction and comparison properties needed from that solution
are verified next.

\begin{proposition}[Model map and boundary data]\label{a:model}
Let $\Phi_*$ be the map of
Proposition~\ref{a:quotient-proposition} for
Equation~\eqref{a:kn-metric}, with axis values
Equation~\eqref{a:axis-values}.  Its ADM energy is $M$, its
angular momentum in the present convention is $J=-aM$, its
charges are $Q,0$, and its horizon area is $4\pi R^2$.
It is independent of Killing time.  Define
\[
 w=\sinh\sigma\sin\vartheta,
 \qquad q_* =\tfrac12\log(W\sin^2\vartheta).
\]
With Euclidean strip derivatives, it satisfies
\begin{equation}\label{a:model-map}
 \divg_G(w\,\dd\Phi_*)=0,\qquad
 -\Delta q_*=|\partial\Phi_*|_G^2,
 \qquad q_*(0,\vartheta)=\log(R^2\sin\vartheta).
\end{equation}
The canonical slice and its tensor and fields extend smoothly
to the bifurcation sphere, including the ordinary axis poles.
Its potentials have the regularity of
Lemma~\ref{a:axis-regularity}; their renormalized components and
$u_*-\log\sin\vartheta$ have bounded derivatives on compact
closed strip ranges.
\end{proposition}

\begin{proof}
The axial radius and quotient are
\begin{equation}\label{a:kn-quotient}
 X_*^2=(W/P)\sin^2\vartheta,
 \qquad
 \gamma_*=-b^2\sinh^2\sigma\sin^2\vartheta\,\dd t^2
                +W\sin^2\vartheta
                    (\dd\sigma^2+\dd\vartheta^2).
\end{equation}
Both Maxwell contractions in Equation~\eqref{a:quotient} have
zero time component.  The same holds for $*_H\dd\mathbf A$,
since $\mathbf A=-a(2Mr-Q^2)\dd t/W$ has curvature of the
form $\dd t\wedge$ a spatial one-form.  Thus the map is
time independent.

An explicit check of all potential constants is useful.  With
$x=\cos\vartheta$, they are
\begin{align}\label{a:kn-potentials}
 \psi_*&=-\frac{Q(r^2+a^2)x}{P},&
 \chi_*&=-\frac{Qar(1-x^2)}P,\notag\\
 z_*&=aMx(3-x^2)
       +\frac{a^3Mx(1-x^2)^2-aQ^2rx(1-x^2)}P.&&
\end{align}
These expressions satisfy the contractions and corrected twist
in Equation~\eqref{a:quotient}, by differentiation.  In
particular the axis values are $(-Q,0,2aM)$ and
$(Q,0,-2aM)$.  To check the angular sign directly, the shift is
\[
 \beta^\phi=-2aM/r^3+O(r^{-4}).
\]
Stationarity and the positive convention $K=\tfrac12\mathcal L_ng$
give $K=-\mathcal L_\beta g/(2N)$ on the canonical slice.
Hence
\[
 K_{r\phi}=-3aM\sin^2\vartheta/r^2+O(r^{-3}),
 \qquad
 \frac1{8\pi}\int_{S_r}K(\nu,\eta)\,\dd A\longrightarrow-aM.
\]
The leading term $Q\dd t/r$ of $\mathfrak a_*$ gives
$F_{tr}=Q/r^2+O(r^{-3})$ and therefore positive electric charge
$Q$ with $E=\iota_nF$.  The magnetic flux is zero.

The horizon is $r=M+b$, or $\sigma=0$.  There
$W=(r^2+a^2)^2=R^4$ by
Equation~\eqref{a:parameter-identities}; its area density is
$R^2\sin\vartheta\,\dd\vartheta\dd\phi$, proving the area
and the boundary value of $q_*$.  In the Euclidean end radius
$\varrho=(b/2)e^\sigma$, one has
$r=\varrho+M+b^2/(4\varrho)$.  Expanding the spatial metric
in the associated Cartesian chart gives
\[
 (g_*)_{ij}=(1+2M/\varrho)\delta_{ij}+O_4(\varrho^{-2}),
\]
and hence ADM energy $M$.  The shift and potential give
$K_*=O_3(\varrho^{-3})$, the electric Coulomb leading term,
and $B_*=O_3(\varrho^{-3})$, with uniform polar derivatives.

Put $\rho_*=bw$.  Equation~\eqref{a:kn-quotient} is static,
with spatial metric $e^{2q_*}(\dd\sigma^2+\dd\vartheta^2)$.
The wave-map equation in Equation~\eqref{a:wave} is therefore
$\divg_G(\rho_*\dd\Phi_*)=0$, giving the first equation in
Equation~\eqref{a:model-map}.  Its spatial Ricci trace is
\[
 -2e^{-2q_*}\Delta q_*
      -\rho_*^{-1}e^{-2q_*}\Delta\rho_*
          =2e^{-2q_*}|\partial\Phi_*|_G^2.
\]
Since $\Delta\rho_*=0$, this proves the second equation.

Finally all the displayed spatial coefficients and potentials
are smooth even functions of $\sigma$ at zero, with the
ordinary polar parity in $\vartheta$.  The apparent lapse
division in the canonical $K_*$ and fields is removable.
Indeed the angular velocity
$\omega_*=a(2Mr-Q^2)/W$ is constant on the horizon:
$\omega_H=a/R^2$.  Thus its $\vartheta$ derivative is
$O(\sigma^2)$ and its $\sigma$ derivative is $O(\sigma)$,
while the lapse is $\sigma$ times a smooth positive even
function.  The electrostatic potential on
$\partial_t+\omega_H\partial_\phi$ is also constant on the
horizon; its nonconstant remainder is $O(\sigma^2)$.
These orders give smooth canonical tensor and field components
after division by the lapse.  The metric has no conical defect
at the axes, and the explicit potentials give precisely the
order-two and corrected order-four axis behavior established
earlier.  For a two-sided extension across $\sigma=0$, the
stationary lapse is taken with its smooth sign,
$b\sinh\sigma\sqrt{P/W}$; the stationary Killing field changes
time orientation across the bridge.  This supplies a smooth
future normal and tensor, whereas using the absolute value of
that lapse would only reflect the chosen stationary time
orientation.  This proves all the regularity assertions.
\end{proof}
